% BEGIN GENERATED ARTIFACT: prf:logical-equivalence-equivalence-relation
\newpage
\phantomsection
\label{prf:logical-equivalence-equivalence-relation}
\LRAProofFor{thm:logical-equivalence-equivalence-relation}

\begin{remark*}[Return]
\hyperref[thm:logical-equivalence-equivalence-relation]{Return to Theorem (Logical Equivalence is an Equivalence Relation).}
\end{remark*}

\begin{theorem*}[Logical Equivalence is an Equivalence Relation]
Logical equivalence is reflexive, symmetric, and transitive on \(\WFF\):
\[
\varphi\equiv\varphi,
\qquad
\varphi\equiv\psi\Longrightarrow\psi\equiv\varphi,
\]
and
\[
(\varphi\equiv\psi\land\psi\equiv\chi)
\Longrightarrow
\varphi\equiv\chi.
\]
\end{theorem*}

\begin{proof}
\textbf{Professional Standard Proof.}~
TODO: supply the compact proof.
\end{proof}

\begin{proof}
\textbf{Detailed Learning Proof.}~
TODO: supply the detailed learning proof.
\end{proof}

\begin{remark*}[Proof structure]
\end{remark*}

\begin{dependencies}
\begin{itemize}
  \item \hyperref[def:logical-equivalence-propositional-logic]{Logical Equivalence}
\end{itemize}
\end{dependencies}

\clearpage
% END GENERATED ARTIFACT: prf:logical-equivalence-equivalence-relation
